% from ega4-2-fr.tex
\oldpage[IV-2]{126}

et du fait que $\operatorname{Ass}(\mathcal O_{Y_\alpha})$ est réduit au point
générique de $Y_\alpha$) entraîne que
$\mathcal H^0_{Y_\alpha/Z_\alpha}(\mathcal O_{Y_\alpha})$ est un
$\mathcal O_{Y_\alpha}$-Module cohérent, et b) découle alors de (5.11.1).

Pour voir que c) entraîne d), on raisonne comme ci-dessus en appliquant c) au
cas où $\mathcal F=(A/\mathfrak q)^\sim$ et $Z=V(\mathfrak J)$;
inversement, on prouve que d) entraîne c) en utilisant encore l'équivalence de
a) et b) dans (5.11.1). Il est évident que c) est un cas particulier de b).
Enfin, le fait que a) (et par suite chacune des autres conditions) est vérifiée
lorsque $A$ vérifie l'une des hypothèses (i), (ii) résulte de (5.11.2), vu la
définition des anneaux universellement caténaires et des anneaux
universellement japonais.

\begin{remarks}[5.11.7]
\label{IV.2.s5.11.7-fr}
\begin{enumerate}
\item[(i)] On ignore si, dans (5.11.5) la condition d) implique a) sans
hypothèse supplémentaire sur $X$; nous verrons plus loin (7.2.4) qu'il en est
bien ainsi lorsque $X$ est un schéma local. De même, nous prouverons que
lorsque $A$ est un anneau local noethérien les quatre propriétés a), b), c),
d) de (5.11.6) sont équivalentes (7.2.4). On ignore si ce résultat s'étend à
tous les anneaux noethériens.
\item[(ii)] Si $A$ vérifie la propriété a) de (5.11.6), il en est de même de
tout anneau de fractions $S^{-1}A$ et de toute $A$-algèbre finie $C$. En effet
tout anneau quotient intègre de $S^{-1}A$ est de la forme
$S^{-1}(A/\mathfrak q)$, où $\mathfrak q$ est un idéal premier de $A$ ne
rencontrant pas $S$; et d'autre part, si $\mathfrak r$ est un idéal premier de
$C$, $\mathfrak p$ son image réciproque dans $A$, $C/\mathfrak r$ est une
$(A/\mathfrak p)$-algèbre finie intègre contenant $A/\mathfrak p$; nos
assertions sont donc des conséquences de (5.10.17, (iii) et (iv)).
\end{enumerate}
\end{remarks}

\subsection{\texorpdfstring{Relations entre les propriétés d'un anneau local noethérien $A$ et
d'un anneau quotient $A/tA$}{Relations entre les propriétés d'un anneau local noethérien A et d'un anneau quotient A/tA}}

On a déjà vu dans (3.4) des relations entre les propriétés de $A$ et de $A/tA$
concernant les idéaux premiers associés, ainsi que les propriétés d'être
intègre ou réduit. On donne dans cette section d'autres relations entre les
propriétés de ces anneaux, liées aux notions de dimension et de profondeur.

\begin{proposition}[5.12.1]
\label{IV.2.s5.12.1-fr}
Soient $A$ un anneau local noethérien, $X=\operatorname{Spec}(A)$, $t$ un
élément de $A$ faisant partie d'un système de paramètres (0, 16.3.6), $X_0$ le
sous-espace fermé $V(t)$ de $X$, $X_1$ l'ouvert complémentaire $X-X_0$. Soit
$Z$ une partie de $X$ telle que toute spécialisation d'un point de $Z$
appartienne à $Z$. On suppose que $X$ est biéquidimensionnel (ce qui sera le
cas si $A$ est équidimensionnel et quotient d'un anneau local régulier (0,
16.5.12)). Alors si l'on pose $Z_0=Z\cap X_0$, $Z_1=Z\cap X_1$, on a
\begin{equation}
\operatorname{codim}(Z_0,X_0)\leq\operatorname{codim}(Z,X)
\leq\operatorname{codim}(Z_1,X_1).
\tag{5.12.1.1}\label{IV.2.s5.12.1.1-fr}
\end{equation}
\end{proposition}

La seconde inégalité résulte de la définition (5.1.3); pour démontrer la
première, on peut se borner à la démontrer lorsque $Z$ est fermé: en effet, par
hypothèse $Z$ est réunion de parties fermées $Z_\alpha$, et si l'on a
$\operatorname{codim}(X_0\cap Z_\alpha,X_0)
\leq\operatorname{codim}(Z_\alpha,X)$ pour tout $\alpha$, on aura aussi
\[
\operatorname{codim}(Z_0,X_0)
=\inf_\alpha\operatorname{codim}(X_0\cap Z_\alpha,X_0)
\leq\inf_\alpha\operatorname{codim}(Z_\alpha,X)
=\operatorname{codim}(Z,X).
\]
Supposons donc $Z$ fermé, de sorte que l'on a
\[
\operatorname{codim}(Z,X)=\dim(X)-\dim(Z)
\]

